\section{Reproducible Evaluation Protocol}
\label{sec:evaluation}

The prespecified evaluation tests canonical-profile compliance, required
reconstruction strength, and fail-closed handling of unsupported, ambiguous,
malformed, over-budget, and incomplete artifacts.

\subsection{Frozen inputs and execution order}
\label{subsec:frozen-run}

A run identifier binds source state, dependencies, target features, policy and
registries, corpus and format digests, external codecs, and server environment.
Any change creates a new append-only record. Failed runs are retained, reruns
cite their predecessors, and pooling requires a prior comparison protocol.

The server follows the gated sequence in
\cref{fig:evaluation-pipeline}. Commands, environment, status, times, and
digests enter the run record. Schema-invalid, incomplete, or hash-inconsistent
reports fail their phase.

\begin{figure*}[t]
\centering
\resizebox{0.94\textwidth}{!}{%
\begin{tikzpicture}[
  phase/.style={draw,rounded corners=2pt,align=center,text width=38mm,
    minimum height=15mm,font=\small,inner sep=3pt},
  freeze/.style={phase,fill=blue!7},
  functional/.style={phase,fill=green!8},
  adversarial/.style={phase,fill=yellow!10},
  terminal/.style={phase,fill=gray!12},
  failbox/.style={draw,rounded corners=2pt,align=center,text width=27mm,
    minimum height=13mm,font=\small,inner sep=3pt,fill=gray!12},
  arrow/.style={-{Latex[length=2mm]},thick},
  fail/.style={-{Latex[length=2mm]},dashed,thick}
]
\node[freeze] (p1) at (0,0) {\textbf{1. Freeze}\\run ID and all input hashes};
\node[freeze] (p2) at (5.2,0) {\textbf{2. Build}\\release build, tests, lint, audit};
\node[functional] (p3) at (10.4,0) {\textbf{3. Functional evidence}\\corpus, ablation, differential, fidelity};
\node[adversarial] (p4) at (10.4,-2.2) {\textbf{4. Adversarial evidence}\\fuzz budgets and injected faults};
\node[adversarial] (p5) at (5.2,-2.2) {\textbf{5. Measurement}\\pinned performance sessions};
\node[terminal] (p6) at (0,-2.2) {\textbf{6. Reduction}\\regenerate artifacts and evaluate $G_{\mathrm{sci}}$};
\node[failbox] (failed) at (15.1,-1.1) {retain failed run\\no controlled\\measurement or release};
\draw[arrow] (p1) -- (p2);
\draw[arrow] (p2) -- (p3);
\draw[arrow] (p3) -- (p4);
\draw[arrow] (p4) -- (p5);
\draw[arrow] (p5) -- (p6);
\draw[fail] (p2.east) -- node[above,sloped,font=\scriptsize]{$G_j=0$} (failed.north west);
\draw[fail] (p3.east) -- (failed.north west);
\draw[fail] (p4.east) -- (failed.south west);
\draw[fail] (p5.east) -- (failed.south west);
\end{tikzpicture}%
}
\caption{Controlled evaluation sequence. Functional or adversarial failure
retains reportable evidence but prevents entry into controlled performance
measurement and release-candidate eligibility.}
\label{fig:evaluation-pipeline}
\end{figure*}

For binary phase-entry and gate variables $E_j,G_j\in\{0,1\}$, the stopping
rule is
\begin{equation}
 E_1=1,\qquad E_{j+1}=E_jG_j,\quad j=1,\ldots,5.
 \label{eq:evaluation-stopping}
\end{equation}
Thus no later measurement compensates for an earlier failed gate. The run
record identifies hardware, operating system, toolchains, CPU affinity,
observed governor and turbo state, background load, and external-tool
identities. A controlled performance run must additionally pin the controls
that the platform permits. Missing required evidence makes
\(G_{\mathrm{sci}}\) false.

\subsection{Property suite}
\label{subsec:property-suite}

The eight invariants in \cref{tab:method-invariants} map one to one to the
distinct tests in \cref{tab:property-suite}. Because the global gate is a
conjunction, aggregate success cannot hide one invalid accepted output or one
fault case that publishes bytes.

\begin{table*}[t]
\centering
\footnotesize
\setlength{\tabcolsep}{4pt}
\caption{Property suite, test oracle, and admissible evidence.}
\label{tab:property-suite}
\begin{tabularx}{\textwidth}{L{0.07\textwidth}L{0.19\textwidth}L{0.39\textwidth}Y}
\toprule
ID & Formal invariant and empirical oracle & Oracle & Evidence unit \\
\midrule
F1 & Complete consumption & Accepted output parses from offset zero and consumed octets equal output length. An unconsumed suffix prevents acceptance & Accepted output and suffix-mutated pair \\
F2 & Component closure & Every parsed component is permitted at its exact grammar location. Unknown or forbidden components are removed through valid reconstruction or blocked & Parsed component path \\
F3 & Output evidence convergence & Generated name, declared output type, byte recognition, and authorizing report resolve to one output profile & Accepted candidate record \\
F4 & Semantic source-independence & A strict semantic path serializes output from decoded media rather than retained source payload & Source-to-candidate provenance trace \\
F5 & Second-pass stability & Under the same policy and profile, a second pass is exact or satisfies the registered semantic closure & First- and second-pass pair \\
F6 & Resource boundedness & Each declared counter, cancellation condition, and deadline returns a typed failure when its bound is exceeded & Boundary and injected-failure case \\
F7 & Validation separate from construction & Every clean candidate is reparsed by a registered authorizing parser distinct from the constructing handler & Accepted candidate and internal authorizing report \\
F8 & Fail-closed result construction and publication & Every failed conjunct, non-clean verdict, error, race, and persistence failure exposes no deliverable candidate and preserves a pre-existing target & Interface and filesystem trace \\
\bottomrule
\end{tabularx}
\end{table*}

F1 appends zero, foreign-object, and malformed suffixes after the terminal
extent. F2 checks parsed component paths, F3 compares normalized profiles, and
F4 verifies that strict semantic output is serialized from decoded media.
Unlike F1--F3, F4 is established by the construction path rather than by an
artifact-observable predicate: byte-level source independence across a
decode--re-encode step cannot be read off the candidate \(y\) alone, so F4 rests
on \cref{asm:semantic-construction} and the code path that serializes only from
the decoded intermediate representation.

Second-pass stability operationalizes F5 but is not an acceptance condition.
Let \(\mathsf{Out}_{p,v}(u)=y\) only for a clean accepted candidate, and let
\(u[y]\) submit that output under the same policy and an equivalent fresh
deadline. For profile \(k\), decoded representation \(\phi_k\), registered
metric \(d_k\), and threshold \(\varepsilon_{p,k}\), define
\begin{equation}
 \mathsf{Close}_{p,k}(y,y')=
 \begin{cases}
  [y'=y],&k\in K_{\mathrm{exact}},\\
  [d_k(\phi_k(y'),\phi_k(y))\le\varepsilon_{p,k}],
    &k\in K_{\mathrm{lossy}},
 \end{cases}
 \label{eq:idempotence}
\end{equation}
and require
\[
\begin{aligned}
 \mathsf{Out}_{p,v}(u)=y&\Longrightarrow
 \exists y'=\mathsf{Out}_{p,v}(u[y]):\\
 &\mathsf{profile}(y')=k\land
 \mathsf{Close}_{p,k}(y,y').
\end{aligned}
\]
Lossy closure is one-step and metric-specific; it implies neither perceptual
equivalence nor bounded accumulated drift. Unlike
\cref{eq:validation-monotonicity}, it allows new bytes and parser reports.

\subsection{Paired mechanism ablation}
\label{subsec:mechanism-ablation}

The paired study measures decision sensitivity within the frozen corpus, not
antivirus accuracy or a deployment-population effect. Let \(P_0\) be the full
strict policy and \(\mathcal{C}\) the conditions in
\cref{tab:mechanism-ablation}. Before ablation, every input is replayed under
\(P_0\). Any mismatch in action, verdict, reason, output digest,
reconstruction, or profile aborts the study.

\begin{table*}[t]
\centering
\scriptsize
\setlength{\tabcolsep}{4pt}
\begin{tabularx}{\textwidth}{L{0.21\textwidth}L{0.28\textwidth}L{0.19\textwidth}Y}
\toprule
Condition & Single registered change from \(P_0\) & Applicable records &
Interpretive endpoint \\
\midrule
\shortstack[l]{\texttt{without-input-}\\\texttt{signature-}\\
\texttt{prescan}} &
Replace input-and-output signature scanning with output-only scanning &
All records &
Observed paired transition after removing pre-reconstruction signature evidence \\
\shortstack[l]{\texttt{without-}\\\texttt{supplied-type-}\\
\texttt{evidence}} &
For records with a supplied type, omit that value and disable its prerequisite.
For records without one, change only the prerequisite &
Two prespecified groups analyzed separately. No pooled intervention estimate &
Observed paired transition within each distinct intervention group \\
\shortstack[l]{\texttt{lowered-}\\\texttt{reconstruction-}\\
\texttt{requirement-}\\\texttt{where-}\\\texttt{configurable}} &
Lower the animated-image, video, and audio requirement from required semantic
reconstruction to structural reconstruction &
Manifest profiles eligible for one of these configurable requirements &
Decision transition under a lower registered reconstruction floor \\
\bottomrule
\end{tabularx}
\caption{Prespecified paired mechanism-ablation conditions. Condition
identifiers match the machine-readable protocol.}
\label{tab:mechanism-ablation}
\end{table*}

Let \(\mathcal I\) be the full corpus and \(I_c\subseteq\mathcal I\) the
preregistered applicable set for condition \(c\). The evidence stream retains
one auditable row for every pair in \(\mathcal C\times\mathcal I\), including an
explicit inapplicable row when \(i\notin I_c\). Estimands use only \(I_c\).
\(E_i^{(c)}\in\{0,1\}\) indicates completion without a pipeline error, and
\(Y_i^{(c)}\in\{0,1\}\) is the observed decision for record \(i\), where one
denotes acceptance and zero denotes blocking. The superscript zero denotes the
eligible full-policy record. The paired transition counts and full-policy
denominators are
\begin{equation}
 \begin{aligned}
 N_{ab}^{(c)}
 &=\sum_{i\in I_c}
   \mathbb{1}\!\left[E_i^{(c)}=1
   \land Y_i^{(0)}=a\land Y_i^{(c)}=b\right],\\
 B_a^{(c)}
 &=\sum_{i\in I_c}\mathbb{1}\!\left[Y_i^{(0)}=a\right],
 \qquad a,b\in\{0,1\}.
 \end{aligned}
 \label{eq:ablation-transitions}
\end{equation}
The two directional summaries are
\begin{equation}
 q_{1\rightarrow0}^{(c)}
 =\frac{N_{10}^{(c)}}{B_1^{(c)}},
 \qquad
 q_{0\rightarrow1}^{(c)}
 =\frac{N_{01}^{(c)}}{B_0^{(c)}}.
 \label{eq:ablation-proportions}
\end{equation}
These are the proportions of full-policy accepts becoming blocks and blocks
becoming accepts. Empty denominators are not estimable. Pipeline errors remain
outside the \(2\times2\) matrix and fail the execution gate.

Let \(m_{ci}\) be the multiplicity of pair \((c,i)\) in the emitted stream.
For condition \(c\), let \(H_c\) contain its prespecified intervention groups:
one group for signature removal, supplied-type-present and supplied-type-absent
groups for type ablation, and animated-image, audio, and video groups for the
reconstruction-floor ablation. Let \(n_{c,h}\) be the applicable count in group
\(h\in H_c\), and let \(N_{\mathrm{err}}\) be the pipeline-error count. The predicates
\(G_{\mathrm{strict}}\), \(G_{\mathrm{context}}\),
\(G_{\mathrm{replay}}\), and \(G_{\Delta P}\) verify, respectively, the strict
starting settings, preservation of the frozen input and context except for the
registered evidence removal, equality of the full-policy replay, and isolation
of the condition-specific policy change. Ablation completeness is
\begin{equation}
 \begin{aligned}
 G_{\mathrm{abl}}={}&G_{\mathrm{strict}}\land G_{\mathrm{context}}
 \land G_{\mathrm{replay}}\land G_{\Delta P}\\
 &{}\land
 \bigwedge_{(c,i)\in\mathcal{C}\times\mathcal{I}}[m_{ci}=1]
 \land \bigwedge_{c\in\mathcal{C}}\bigwedge_{h\in H_c}[n_{c,h}\geq1]
 \land [N_{\mathrm{err}}=0].
 \end{aligned}
 \label{eq:ablation-gate}
\end{equation}
The exact-pair term requires a complete \(3|\mathcal{I}|\) audit matrix without
duplicates or omissions. Inapplicable rows do not enter
\cref{eq:ablation-transitions}. The proportions are descriptive frozen-corpus
properties without a population interval, significance test, causal
interpretation, malware-detection rate, or prevalence claim.

\subsection{Differential output validation}
\label{subsec:differential-validation}

Every accepted output is checked by the authorizing parser and at least two
applicable external format or profile validators. Distinct families use
different executable or library lineages. Separate process boundaries alone do
not establish diversity. Reports bind output digests to executable identities,
detected profiles, and outcomes. Any missing, rejecting, or disagreeing result
fails the gate. These external checks corroborate decodability or profile
recognition only. They do not independently establish complete consumption,
component closure, the expected action, or the fixture-derived decision oracle,
and they cannot exclude common-mode errors.

\subsection{Coverage-guided fuzzing and bounded checks}
\label{subsec:fuzzing}

Fuzzing records fixed budgets, seeds, repetitions, sanitizers, corpora, and
failure handling, following established evaluation guidance~\cite{klees2018}.
Targets cover top-level dispatch, complete output validation, and native
animated-image, MP4, WebM, Ogg, and audio paths. Structured seeds preserve outer
headers while mutation reaches lengths, offsets, checksums, lacing, and
sequences. Unique failures are minimized, hashed, reproduced, and added to S7
when appropriate.

Budget exhaustion with no discovered failure means only ``no failure observed''
under that campaign. Native-interface and policy-deserialization properties
receive separate regression and static review, not dedicated fuzz-coverage claims.
Bounded checks cover named finite arithmetic and parser domains. Neither
fuzzing nor static review verifies an external codec.

\subsection{Crash, cancellation, and publication campaigns}
\label{subsec:fault-campaigns}

The fault schedule exercises cancellation, deadlines, child launch and timeout,
bounded diagnostics, output-size termination, validation failure, caught panic,
and attempted non-clean publication. The executed cases record the injection point, terminal
state, paths, child status, and cleanup against
\begin{equation}
 \mathrm{FailureBeforePublication} \Longrightarrow
 \mathrm{DeliverableBytes}=0.
 \label{eq:no-publication}
\end{equation}
The destination is monitored and any pre-existing object is hashed. Allocator,
temporary-file, and persistence failures require explicit injection evidence
before \(G_{\mathrm{fault}}\) can pass. Code review can motivate the expected
behavior but cannot substitute for an executed fault case.

\subsection{Performance protocol}
\label{subsec:performance-protocol}

After functional gates pass, ten sessions invoke the public byte-defense API
over corpus permutations generated from a recorded seed for each session. The matrix stratifies
format, reconstruction level, input and decoded size, media structure, carrier,
and outcome. Cold and warm records retain order, time, profile,
reconstruction, outcome, and errors. Blocked and timed-out calls remain in
their path denominator. Here, ``cold'' denotes the first timed invocation after
an already constructed defender. Policy normalization and signature-engine
construction are outside the timed interval.

Latency summaries report the median, 95th percentile (p95), and 99th percentile
(p99). Sequential single-worker service rate is not concurrent throughput.
Hashed resource logs provide peak resident set size (RSS). Output expansion for
clean deliverable \(i\) is
\begin{equation}
 e_i=\frac{|y_i|}{|x_i|},
 \label{eq:expansion}
\end{equation}
with zero-length input retained as a separate invalid or blocked group.

\subsection{Fidelity and usability}
\label{subsec:fidelity-protocol}

Security acceptance and fidelity remain separate. Every accepted output is
evaluated under a frozen fidelity policy. Lossless profiles compare the
prespecified decoded dimensions, channels, frames, samples, and timing fields.
Structural profiles test the exact properties named by their profile. Lossy
paths use prespecified metrics and drift bounds. Tool error, missing or nonfinite
metrics, threshold failure, or an unevaluated output fails the fidelity gate.
Passing fidelity never enlarges the security claim.

\subsection{Statistical analysis}
\label{subsec:statistical-analysis}

Corpus endpoints report \(x/n\) under the denominators in
\cref{subsec:corpus-denominators}. Because purposive mutations are clustered by
base artifact, no binomial population interval is attached to decision,
differential, fidelity, or invariant outcomes. Zero denominators are not
estimable. The first release has no confirmatory policy or version hypothesis.
Its paired discordances are descriptive.

Median latency uses 10,000 hierarchical bootstrap replicates over sessions,
dependence clusters, and base artifacts~\cite{efron1979}. The resulting
percentile 95\% interval quantifies variation within the controlled design, not
deployment-population uncertainty. Expansion is reported as an arithmetic mean
without a bootstrap interval. Cold and warm paths remain separate. Exploratory
subgroups are labeled, missing observations are not imputed, and every output
table traces to its frozen source digests.

\subsection{Release gates and claim authorization}
\label{subsec:release-gates}

Let each gate take a three-valued state
\(\mathbb G=\{\mathsf P,\mathsf F,\mathsf B\}\) for pass, fail, and blocked.
For a finite gate vector \(\mathbf g\), define
\begin{equation}
 \bigwedge\nolimits^\star\mathbf g=
 \begin{cases}
  \mathsf F,&\exists j:g_j=\mathsf F,\\
  \mathsf P,&\forall j:g_j=\mathsf P,\\
  \mathsf B,&\text{otherwise}.
 \end{cases}
 \label{eq:three-valued-gate}
\end{equation}
The scientific gate is
\begin{equation}
\begin{aligned}
 G_{\mathrm{sci}}=\bigwedge\nolimits^\star\bigl(&
 g_{\mathrm{prov}},g_{\mathrm{schema}},
 (g_r)_{r\in\mathcal C},(g_{\mathrm F j})_{j=1}^{8},
 g_{\mathrm{abl}},g_{\mathrm{diff}},g_{\mathrm{fidelity}},\\
 &g_{\mathrm{fuzz}},g_{\mathrm{fault}},g_{\mathrm{perf}},
 g_{\mathrm{regen}}\bigr).
\end{aligned}
 \label{eq:scientific-gate}
\end{equation}
Here \(\mathcal C\) is the frozen corpus; provenance freezes all input
identities, schema validates cross-file consistency, \(g_r\) binds each
record-level oracle and evidence, and \(g_{\mathrm F j}\) instantiates
\cref{tab:property-suite}. The remaining gates require their complete
prespecified ablation, external-validation, fidelity, fuzz, fault, performance,
and regeneration schedules. Controlled or release claims require
\(G_{\mathrm{sci}}=\mathsf P\) and
\texttt{release\_eligibility.eligible=true}; \(\mathsf F\) or \(\mathsf B\)
retains finite observations but forbids those claims.
